\section{创新性、实用性与安全设计}
\label{sec:conclusion}

\subsection{主要创新点}

\textbf{第一，系统级闭环而非单点模型。}本作品将欺诈样本注入、交易流回放、Flink实时计算、模型API、Dashboard展示、人工反馈和重训接口串联为闭环，能够直接演示“新增交易进入系统、实时评分、告警输出、人工复核、模型迭代”的完整过程，更贴近真实金融风控系统。

\textbf{第二，窗口特征与画像特征在线一致。}离线训练通过\code{add\_offline\_window\_features}构造窗口统计，在线Flink作业使用同一套\code{compute\_window\_feature\_dict}逻辑实时计算，降低训练/推理特征偏差。画像数据通过Redis按实体ID补齐，兼顾实时性和特征丰富度。

\textbf{第三，可解释图挖掘与GraphSAGE旁路并行。}系统围绕共享设备、IP、商户和收款方生成团伙证据链，同时以GraphSAGE旁路实验验证图嵌入价值。线上决策不直接使用未经充分验证的图嵌入，体现对创新探索与主链路稳定性的平衡。

\textbf{第四，阈值沙盒连接模型指标与业务研判。}模型分数被转换为风险等级、处置动作和理由码，阈值沙盒允许实时观察不同阈值下的精度、召回、误拦、漏过和审核工作量，使模型优化能够直接映射到业务成本。

\textbf{第五，微批推理提升工程可用性。}在线评分支持单笔与微批两种模式。微批模式通过批大小和等待时间控制延迟上界，在高并发流量下显著减少API调用次数，为后续横向扩展和生产化部署预留空间。

\textbf{第六，反馈驱动但受控发布。}系统采用candidate/latest/rollback三段式模型生命周期。反馈样本先训练候选模型，经指标核验和人工确认后发布，异常时可快速回滚，避免“自动学习”直接导致线上模型退化。

\subsection{实用性与业务适配}

本系统面向第三方支付交易风控，输出内容与业务处置流程一致：低风险交易放行，中风险交易人工复核，高风险交易拒绝并写入告警。Kafka主题天然适合接入上游支付网关和下游告警、审核、风控策略系统；FastAPI模型服务可独立扩容；Redis画像服务可替换为企业内部画像库或特征平台；Dashboard可作为比赛演示版，也可扩展为运营监控台。

在业务规则层面，系统覆盖常见反诈场景：异常大额、短时高频、多账户共设备、代理IP、风险商户集中、欺诈邻域和团伙命中。上述特征均可映射到真实金融机构风控经验，便于与规则引擎、黑名单、设备指纹系统和人工审核流程对接。

\subsection{成熟度与可扩展性}

代码结构按功能模块拆分，\code{fraudsim/models}负责模型适配，\code{fraudsim/training}负责训练，\code{fraudsim/streaming}负责流式链路，\code{fraudsim/api}负责服务接口，\code{fraudsim/dashboard}负责前端展示，\code{scripts}负责数据、启动和压测。新增模型只需实现统一适配器并注册；新增欺诈场景可在注入脚本中扩展；新增画像字段可通过特征配置自动进入训练与推理；新增下游系统可订阅Kafka输出主题。

系统提供Dockerfile、Dockerfile.flink、Dockerfile.graph和docker-compose.yml，使运行环境可复制；同时提供边缘实时侧与中心训练侧Kubernetes清单。边缘清单包含双副本Model API、Flink、Kafka、Redis、Secret和NetworkPolicy，中心清单包含候选模型训练与GraphSAGE实验CronJob。离线结构校验结果为20个资源有效、0个无效、0个错误。

\subsection{安全与合规设计}

金融反诈系统处理交易、设备、IP和账户行为数据，必须满足全生命周期安全要求。当前作品使用仿真与脱敏数据，并已实现部分基础安全能力：

\begin{enumerate}
  \item \textbf{数据最小化与脱敏：}训练和展示层仅使用实体ID哈希值，不展示姓名、证件号、手机号、银行卡号等直接身份标识；演示数据使用仿真样本；
  \item \textbf{敏感接口认证：}配置\code{FRAUDSIM\_API\_KEY}后，预测、反馈、重载、演示任务、模型发布和回滚等接口要求\code{X-API-Key}；
  \item \textbf{双通道审计：}敏感操作同时写入本地\code{audit.jsonl}和Kafka的\code{audit\_events}主题，Kafka不可用时仍保留本地审计；
  \item \textbf{模型发布保护：}候选模型发布前自动备份当前线上模型，支持快速回滚；
  \item \textbf{部署隔离：}Kubernetes清单使用Secret注入密钥，并通过NetworkPolicy限制Model API访问范围；
  \item \textbf{生产增强方向：}实际部署中还需启用TLS/SASL、OAuth2或mTLS、实体ID令牌化、对象存储加密、RBAC与不可篡改审计留存。
\end{enumerate}

\subsection{局限性与后续工作}

当前系统仍有三点需要继续增强。第一，代码仓库未包含大体积数据和模型产物，最终提交应附带可复现的数据准备说明与关键实验产物。第二，GraphSAGE当前仅完成受控旁路实验，后续需补充时间切分、消融实验和稳定性验证后再评估是否融合线上决策。第三，Kubernetes清单已完成结构校验但未在真实多节点集群部署，生产级环境仍需补齐镜像仓库、StorageClass、TLS/SASL、Ingress、RBAC和故障恢复压测。

总体来看，FP-FraudSim已经覆盖赛题一所要求的主要技术面：能够按参数模拟跨时段、跨渠道与团伙欺诈交易，能够在Kafka + Flink链路中实时检测并输出风险等级，能够通过图证据、理由码与阈值沙盒辅助研判，并通过反馈样本、候选发布、审计和回滚形成受控迭代闭环。
